\section{The split Eisenstein eightfold family: an exact \(G_{34}\) construction}
\label{sec:g34-new}

Put \(K=\mathbb Q(\omega)\), where \(\omega^2+\omega+1=0\), and
\(\mathcal O=\mathbb Z[\omega]\).  We give the branch datum, reconstruct its
parabolic monodromy, prove that the period map has rank exactly \(16\), and
supply the algebraic determinant classes needed for the transfer.  The final
bound is integral and effective.  Initially the target lattice is
\(\mathcal O\)-stable; the last paragraph treats an arbitrary fixed order.

\paragraph{Attribution and status of the inputs.}
The ten roots, eight cocycle representatives, and seventeen intervals below
are numerical data from Appendix A of Benjamin Severn's unrefereed manuscript
\cite{G34:Severn}, Hexagon \(2610.00387v2\), internal version \(1.5\), dated
21 October 2026.  The actual PDF and TeX title is \emph{Weil classes on abelian
sixfolds and eightfolds of Weil type for \(\mathbb Q(\sqrt{-3})\)}.  The portal
heading instead mentions only sixfolds.  The downloaded TeX has SHA256
\texttt{fc15af16599b399d8f111f6f0804916f}\allowbreak\texttt{dd39dc627e83adb5ff4e1788babc75ec},
agreeing with the portal's advertised prefix.  The source log preserves the
metadata discrepancy and the hashes of the actual files read.

Severn's Theorem A claims the split eightfold result; his Theorems B and C
concern sixfolds and powers.  Those theorems remain conditional if invoked as
inputs.  None is invoked here.  In particular, neither his dominance criterion,
his supplier-transfer theorem, nor his printed monodromy or Gram matrices is
a premise.  The numerical witness is checked by a separate exact reconstruction;
the arguments below use the published inputs identified at each step.
Markman and 032 are not used in this section.  The check events distinguish
exact computations, whole-proof checks, and published background.

\subsection*{D-G34-01. Group identification, branch datum, and the eightfold}

Equip \(K^6\) with \(\langle x,y\rangle=\overline{x}^{\,t}y\).  The roots
\(v_i\) are the rows of the following array, regarded as column vectors when
forming matrices:
\[
\begin{array}{c|rrrrrr}
i&\multicolumn{6}{c}{v_i^t}\\\hline
1&1&-1-\omega&1&\omega&-1-\omega&\omega\\
2&1&0&0&1+\omega&0&0\\
3&0&1&1+\omega&0&0&0\\
4&1&\omega&-1-\omega&\omega&-1-\omega&1\\
5&1&-1-\omega&-1-\omega&-1-\omega&-1-\omega&\omega\\
6&1&-1-\omega&\omega&\omega&\omega&\omega\\
7&0&0&1&0&0&-1\\
8&1&\omega&\omega&-1-\omega&1&-1-\omega\\
9&1&-\omega&0&0&0&0\\
10&0&1&0&0&0&-1
\end{array}
\]
Define
\[
 s_i=1-\frac{2v_i\overline v_i^{\,t}}{\langle v_i,v_i\rangle},\qquad
 (t_1,\ldots,t_{20})=(s_1,s_1,\ldots,s_{10},s_{10}).
\]
Each root has norm \(2\) or \(6\), and each \(s_i\) is an involutive
reflection.  Consequently \(t_1\cdots t_{20}=1\).

\textbf{Claim.} The \(s_i\) generate \(G_{34}\), of order \(39\,191\,040\).
Their natural representation is absolutely irreducible over \(K\), has
character field \(K\), and its rational simple block is \(M_6(K)\).

\begin{proof}
Consider the following set \(R\) of norm-two vectors:
\[
\begin{split}
 R_{\rm sp}&=\{\omega^ae_i-\omega^be_j:
          i\ne j,\ a,b\in\{0,1,2\}\},\\
 R_{\rm de}&=\left\{\frac{\pm1}{1+2\omega}
       (\omega^{a_1},\ldots,\omega^{a_6}):
       a_i\in\{0,1,2\},\ \sum a_i\equiv0\pmod3\right\}.
\end{split}
\]
They have cardinalities \(270\) and \(486\), respectively.  The exact
reconstruction checks that every \(s_i\) permutes \(R=R_{\rm sp}\cup R_{\rm de}\),
that their generated action is transitive on its \(756\) elements, and that
the permutation group has order \(39\,191\,040\).  Its stabilizer-chain
orbit lengths are
\[
                   756,\quad80,\quad12,\quad9,\quad6,
\]
whose product is the stated order.  The permutation action is faithful because
\(R\) spans \(K^6\).

Davis \cite[Theorem 5.1 and its proof, pp.~312--313]{G34:Davis} describes
\(G_{34}\) by the reflections in the hyperplanes
\(x_i-x_j=0\), \(x_1-\omega x_2=0\), and
\(x_1+\cdots+x_6=0\).  He gives the invariant degrees
\(6,12,18,24,30,42\), whose product is \(39\,191\,040\), and the same
\(270+486\) root set.  One root reflection in our group, together with
transitivity on \(R\), gives the reflections in every root direction.  In
particular it gives Davis's standard generators.  Thus the published
\(G_{34}\) is a subgroup of the computed group; equality of the exact orders
proves equality of the groups.  This is an identification by containment and
order, not just by a remembered order and dimension.

The roots have rank six and their nonorthogonality graph is connected, both
checked over \(K\).  If a complex invariant subspace \(W\) is given, then
\(W^\perp\) is invariant.  Since \(1-s_i\) has rank one, its root line
lies in \(W\) or \(W^\perp\).  Nonorthogonal roots cannot belong to opposite
summands.  Connectedness puts all roots in the same summand, and their span
is all of \(\mathbb C^6\).  This proves absolute irreducibility.

The exact permutation calculation also verifies that \(\omega I\) belongs
to the group.  The character values lie in \(K\), while
\(\operatorname{tr}(\omega I)=6\omega\) generates \(K\) over \(\mathbb Q\).
The character field is therefore exactly \(K\).  The displayed realization
over that field has Schur index one, giving the block \(M_6(K)\).
\end{proof}

Riemann existence now supplies a connected \(G_{34}\)-cover
\(C\to\mathbb P^1\) with this product-one generating tuple.  The corresponding
Hurwitz component is a finite cover of the ordered configuration space
\(M_{0,20}\), hence has dimension \(17\).  Its genus is
\[
 2g(C)-2=|G_{34}|\left(-2+20\left(1-\frac12\right)\right)
        =8|G_{34}|,
 \qquad g(C)=156\,764\,161.
\]
For the natural character \(\sigma\) and its conjugate, every inertia
eigenvalue list is \((1,1,1,1,1,-1)\).  Chevalley--Weil in the precise form
of \cite[Theorem 2.10]{G34:FGP} gives holomorphic multiplicity
\[
                         -6+20\cdot\frac12=4
\]
for each embedding.  A rational primitive idempotent of \(M_6(K)\) cuts out
an abelian variety \(A\) of dimension eight and Weil signature \((4,4)\).
The parabolic multiplicity space has \(K\)-dimension eight.  The following
calculation verifies the latter dimension independently at the cocycle level.

\subsection*{D-G34-02. The exact parabolic and boundary certificate}

Use left cocycles \(u(gh)=u(g)+g u(h)\).  Write
\[
 u_i=u(\gamma_i)=\lambda_i v_{\lceil i/2\rceil},\qquad
 Q_i=t_1\cdots t_{i-1}.
\]
The parabolic relation and the coboundaries are
\[
 \sum_{i=1}^{20}Q_i\lambda_i v_{\lceil i/2\rceil}=0,
 \qquad u_i=(1-t_i)a.
\]
In \(\lambda\)-coordinates the coboundary matrix has row
\(2\overline v_{\lceil i/2\rceil}^{\,t}/
\langle v_{\lceil i/2\rceil},v_{\lceil i/2\rceil}\rangle\).
The relation matrix and coboundary matrix have ranks six and six, and their
product is zero.  Thus the quotient dimension is \(20-6-6=8\).

For \(j=1,\ldots,4\), let \(b_j\) have
\(\lambda_{2j-1}=\lambda_{2j}=1\) and all other entries zero.  The last four
vectors are supported at the ordered indices
\[
                       (1,3,5,7,10,12,14,16,18,20)
\]
and have coefficient rows
\[
\begin{aligned}
 b_5:;&(1,0,0,0,1+\omega,1,0,-1-\omega,0,-1+\omega),\\
 b_6:;&(0,1,0,0,1+\omega,\tfrac23-\tfrac23\omega,-1,
                     -\tfrac23-\tfrac13\omega,0,1+2\omega),\\
 b_7:;&(0,0,1,0,-\tfrac13-\tfrac23\omega,-\tfrac23-\tfrac13\omega,
                          1+\omega,0,1+\omega,-\omega),\\
 b_8:;&(0,0,0,1,1+\omega,1,-2-\omega,-1-\omega,0,3\omega).
\end{aligned}
\]
All eight satisfy the relation.  Adjoining the six coboundary columns gives
rank fourteen, so these are a quotient basis, rather than just candidate
representatives.  Put \(I=\langle b_1,b_2,b_3,b_4\rangle_K\).

For an auxiliary pairing certificate, define
\[
 F_i(u)=\sum_{j<i}Q_j u_j,\qquad
 S(u,v)=\sum_i\left\langle Q_i^{-1}F_i(u)+\frac12u_i,v_i\right\rangle.
\]
The expression is the customary left-cocycle cup expression, comparable with
\cite[Theorem 2.5]{G34:DWdual}.  The proof below does not need a convention
conversion identifying it with the geometric cup product.  It is used only
as an explicitly verified auxiliary matrix.  Exact evaluation gives
\(S^*=-S\), \(\operatorname{rank}S=8\), zero upper-left block, and
upper-right block \(\operatorname{diag}(3,1,1,3)\).  All entries are saved
in the computation certificate.

The intervals and return powers are as follows:
\[
\begin{array}{c|rrr|r@{\qquad}c|rrr|r}
 j&a&b&k&\operatorname{rank}N_j&j&a&b&k&\operatorname{rank}N_j\\\hline
1&1&8&1&4&10&9&12&1&2\\
2&1&18&1&1&11&9&16&1&4\\
3&2&8&2&3&12&9&17&2&4\\
4&3&8&1&3&13&9&18&1&4\\
5&4&8&2&2&14&11&12&1&1\\
6&5&6&1&1&15&13&14&1&1\\
7&5&8&1&2&16&13&16&1&2\\
8&7&8&1&1&17&15&16&1&1\\
9&9&10&1&1&&&&&
\end{array}
\]
They are pairwise nested or disjoint and each gives a stable partition.
They therefore determine seventeen disjoint nonisotopic essential circles
on the twenty-punctured sphere, a maximal degeneration in
\(\overline M_{0,20}\).  For an interval \([a,b]\), put
\(c=t_a\cdots t_b\).  The listed \(k\) satisfies \(c^k=1\).

The \(k\)-th boundary twist fixes the tuple exactly.  Its pullback is computed
from
\[
\begin{split}
 u(c^k)&=(1+c+\cdots+c^{k-1})
             \sum_{i=a}^b(t_a\cdots t_{i-1})u_i,\\
 u'_i&=u_i+(1-t_i)u(c^k)\quad(a\le i\le b),\qquad
 u'_i=u_i\quad(i\notin[a,b]).
\end{split}
\]
Indeed the twist conjugates each inside loop by the boundary word
\(\gamma_a\cdots\gamma_b\); the cocycle identity gives the formula, and
\(c^k=1\) gives the returning action.  Reducing the resulting cocycles
modulo the verified coboundary space gives matrices
\[
 T_j=1+N_j,\qquad N_j=
 \begin{pmatrix}0&X_j\\0&0\end{pmatrix},\qquad
 N_iN_j=0,\quad T_j^*ST_j=S.
\]
The ranks are in the table.  These matrices have been derived from the roots,
not copied from the printed monodromies.

Let \(D=\operatorname{diag}(-3,-1,-1,-3)\).  The matrices \(-DX_j\) are
Hermitian positive semidefinite: all fifteen nonempty principal minors of
each are nonnegative rational numbers, checked exactly.  The leading principal
minors of their sum are
\[
                         96,\quad2379,\quad58722,\quad2344041.
\]
They are positive.  Every positive weighted sum of the \(N_j\) consequently
has rank four, with kernel and image both \(I\).

Flatten the sixteen matrices with indices
\[
                    1,2,3,4,5,6,7,9,10,11,12,13,14,15,16,17
\]
row by row into vectors of length sixteen over \(K\).  Their determinant is
\[
                                  1728\ne0.
\]
Thus the \(K\)-span of all \(X_j\) is \(M_4(K)\), and their complex span
under either embedding is \(M_4(\mathbb C)\).  The supplementary rational
flattening has rank sixteen and a nonzero minor \(64\); it is the
\(K\)-minor that proves complex tangent rank, so no rational-dimension
substitution is being made.

\paragraph{Reproducible finite certificate.}
The executable \texttt{g34\_reconstruct.py}, invoked with
\texttt{--group-order}, and the complete output
\texttt{g34\_reconstruction\_result.json} accompany this section in its
agent directory.  The latter contains every entry of \(S\) and all seventeen
\(X_j\), the selected minor, every principal minor, the ranks, and the group
orbit lengths.  Entries are encoded as two rational coefficients of
\(1,\omega\).  The input consists only of the ten roots, the eight
representatives, and the interval list, all printed above.  Fraction arithmetic,
Gaussian elimination, determinants, the parabolic quotient, and twist pullback
are implemented directly over \(\mathbb Q[\omega]/(\omega^2+\omega+1)\).
Sympy is used only for the exact finite permutation-group order.  The
additional \texttt{g34\_matrices.tex} prints the output matrices for a reader
who prefers a typeset certificate.  The saved checks are exact computations,
not proof-assistant formalizations.

\subsection*{D-G34-03. The period differential has rank exactly sixteen}

\textbf{Statement.} This Hurwitz family has smooth points where its period
map has rank sixteen.  Its image contains an analytic open subset of the
split signature-\((4,4)\) unitary period domain, and its algebraic moduli map
is dominant onto its connected Weil component.

\begin{proof}
Pass to a finite marking cover so that the \(G\)-cover and its multiplicity
factor are families.  Near the maximal degeneration above, the stable pointed
sphere has seventeen independent plumbing parameters.  On the punctured
polydisc the Hurwitz cover has finite degree.  Since its fundamental group
is \(\mathbb Z^{17}\), a finite Kummer base change kills its finite
permutation monodromy.  Choose the exponents to be multiples of the listed
return powers.  A further common multiple kills any finite central marking
action.  A base point outside the interval circles and a marking of its lifted
fibre give the cocycle pullback calculated above; equivalently this is the
geometric-monodromy identification in
\cite[Lemma 2.2 and Theorem 2.5]{G34:DWmon}.

Every positive plumbing loop is the positive twist about its vanishing
circle.  The passage from a twist to cohomological pullback may reverse all
signs, but does not independently reverse individual signs.  Therefore the
logarithms of the resulting unipotent monodromies are
\[
                        M_j=\varepsilon c_jN_j,
             \qquad c_j\in\mathbb Z_{>0},\quad\varepsilon\in\{1,-1\},
\]
with one common \(\varepsilon\).  Powers only multiply these logarithms,
since \(N_j^2=0\).

Write the punctured coordinates as \(z_j=e^{2\pi i w_j}\).
Schmid's multivariable nilpotent-orbit theorem
\cite[Theorem (4.12), p.~233]{G34:Schmid} gives
\[
              F(w)=\exp\!\left(\sum_j w_jM_j\right)\Psi(z),
\]
where \(\Psi\) extends holomorphically to the polydisc with values in the
compact dual.  Let \(F_0=\Psi(0)\).  Restrict to the positive diagonal
\(z_1=\cdots=z_{17}=z\).  Its logarithm is \(M=\sum_jM_j\).
The semidefinite certificate, which is purely a statement about the computed
matrices and does not require identifying the auxiliary form with the
polarization, gives
\[
              \operatorname{rank}_K M=4,\qquad
                         \operatorname{im}M=\ker M=I.
\]
On the underlying rational \(H^1\), the rank is eight.  The monodromy weight
filtration centered at one is
\[
                 W_0=W_1=I_{\mathbb Q},\qquad W_2=V_{\mathbb Q}.
\]
The diagonal restriction has the same limiting filtration \(F_0\).
Schmid's limiting mixed-Hodge-structure theorem
\cite[Theorem (6.16), p.~255]{G34:Schmid} applies.  Since
\(F_0^0=V_{\mathbb C}\) and \(F_0^2=0\), the weight-zero graded piece is
of type \((0,0)\).  More explicitly, both the filtration and its conjugate
have index-zero term equal to the whole space, so Hodge indices on a pure
graded piece are nonnegative; at weight zero they must both be zero.
Consequently
\[
                     F_0^1\cap I_{\mathbb C}=0.
\]
On either \(K\)-embedding block \(V_\sigma\), dimensions are four and four,
and hence
\[
                         V_\sigma=I_\sigma\oplus F_{0,\sigma}^1.
\]

The tangent space of the compact dual \(\operatorname{Gr}(4,V_\sigma)\)
at \(F_{0,\sigma}^1\) is
\(\operatorname{Hom}(F_{0,\sigma}^1,V_\sigma/F_{0,\sigma}^1)\).
Under the displayed decomposition, the infinitesimal action of \(M_j\)
is its induced map from \(F_{0,\sigma}^1\) to \(I_\sigma\).
Changing from the fixed quotient complement in the finite calculation to
\(F_{0,\sigma}^1\) is a fixed invertible map.  Thus these tangent maps
span the same dimension as the \(c_jX_j\), namely all sixteen dimensions
by the determinant \(1728\).

This limiting rank persists at an interior point.  Differentiate the period
map with respect to \(w_j\), and translate its tangent vectors back by
\(\exp(-\sum w_iM_i)\).  The result is
\[
                   M_j\Psi(z)+2\pi i z_j\frac{\partial\Psi}{\partial z_j}.
\]
The second term tends to zero, and the first to \(M_jF_0\).  In a holomorphic
coordinate chart near \(F_0\), take the determinant of the sixteen columns
selected in the finite certificate.  Its limit is nonzero.  It is therefore
nonzero for some sufficiently small point with all coordinates nonzero.
Translation by the group preserves rank.  The smooth period map consequently
has rank at least sixteen; the unitary period domain has dimension
\(4\cdot4=16\), so the rank is exactly sixteen.

It remains to verify that the component is split using the actual geometric
polarization.  The logarithm \(M_1\) is skew-adjoint for that polarization,
its square is zero, and its image is \(I\).  For any \(u,v\), the pairing
of \(M_1u,M_1v\) is minus the pairing of \(u,M_1^2v\), hence zero.  Thus
\(I\) is an actual \(K\)-rational totally isotropic four-space in the
rank-eight Hermitian multiplicity space.  The form is split and its
determinant class is \(1\).  This reasoning is independent of the auxiliary
cup-form convention.  The period image lies in the required component;
rank sixteen makes it analytically open locally, and its algebraic moduli
map is dominant.  The source-dimension inequality was never used as a
sufficient dominance test.
\end{proof}

\subsection*{D-G34-04. The exceptional-cohomology step in Schoen's construction}

\textbf{Statement.} In the paired-branch blowup construction of
\cite[\S2]{G34:Schoen}, the inverse image of each exceptional divisor has
zero cohomology in every faithful cyclic character.  This holds for every
cyclic order, in particular for inertia exponents \(m/2\).

\begin{proof}
Work over \(\mathbb C\).  Suppose a normalized cyclic \(m\)-cover has
normal-crossing branch divisors, and two divisors of nonzero exponents
\(a,-a\pmod m\) are blown up along their transverse intersection \(Z\).
The exceptional divisor \(p:E\to Z\) is a \(\mathbb P^1\)-bundle.
Its total exponent is zero modulo \(m\), so the normalized pullback is
generically unramified along it.  The two strict branch transforms meet
\(E\) in disjoint sections.  Every other branch divisor restricts to the
full inverse image of its intersection with \(Z\).  These statements can
also be checked in the two blowup charts of the local equation
\[
             y^m=(\text{unit})s^a t^{m-a}\prod_\ell u_\ell^{a_\ell}:
\]
remove the \(m\)-th power of the exceptional parameter.

Let \(\beta:F\to E\) be the finite inverse image and \(\chi\) a faithful
character of \(\mu_m\).  Put
\(\mathscr F_\chi=(\beta_*\mathbb C)_\chi\).
At a point with nontrivial inertia, its stalk is zero.  Indeed the fibre
representation is a permutation representation, and its \(\chi\)-part
is nonzero exactly when \(\chi\) is trivial on the stabilizer.  Away from
the branch locus it is a rank-one local system.

On a \(\mathbb P^1\)-fibre not contained in another branch divisor, its
restriction is therefore extension by zero of a rank-one local system on
\(\mathbb P^1\setminus\{0,\infty\}\) with monodromy
\(e^{2\pi i a/m}\ne1\).  On a fibre contained in another branch divisor,
the restriction is zero.  A rank-one local system on \(\mathbb C^*\) with
monodromy \(\lambda\ne1\) has zero cohomology: its cellular complex has
differential \(\lambda-1\), an isomorphism.  Duality gives zero compactly
supported cohomology as well.  Proper base change therefore gives
\(Rp_*\mathscr F_\chi=0\).  Since \(\beta\) is finite and proper,
\(R\beta_*\mathbb C=\beta_*\mathbb C\), and hence
\[
                            H^*(F,\mathbb C)_\chi=0.
\]
One may apply the exact standard forms SGA4 XII, Theorem 5.1(iii), and
XVI, Theorem 4.1, by first using finite coefficients of residue characteristic
prime to \(m\), adjoining the \(m\)-th roots of unity, and then passing
to characteristic-zero coefficients and classical comparison
\cite{G34:SGA4}.  The cellular differential remains invertible with those
coefficients.

The inverse image of the original centre \(Z\) has the same faithful-character
vanishing: its inertia contains the nontrivial subgroup generated by \(a\),
so the corresponding pushforward eigensheaf is zero at every point.
The blowup is an isomorphism on the common complement.  Comparing the two
compact-support localization sequences, whose closed terms both have zero
faithful-character cohomology, identifies the faithful-character parts
before and after the blowup.  In particular no new primitive character is
introduced.  At subsequent stages the same argument applies, since the
branch transforms retain normal crossings.  This proves the assertion
needed in Schoen's Lemma 2.8.
\end{proof}

The repair is necessary for the proof being used: on p.~18 of
\cite{G34:Schoen}, a prime \(p\mid m\) is chosen such that the degree-\(p\)
quotient remains ramified at a selected point.  For \(m=12\) and local
exponent \(6\), all prime-degree quotients are unramified there.  The direct
faithful-character argument above avoids that choice.  The remaining published
paired-branch construction is used through Theorem 2.0, Lemmas 2.2--2.6,
and Corollary 3.1 of \cite{G34:Schoen}; their hypotheses are verified next.

\subsection*{D-G34-05. Algebraic cyclic suppliers for every needed character}

\textbf{Statement.} Let \(C\to\mathbb P^1\) be a connected finite
\(G\)-cover with all inertia groups of order two.  For a subgroup
\(H'\le G\) and a linear character \(\chi:H'\to\mathbb C^*\) of order
\(m\ge3\), put \(C_\chi=C/\ker\chi\), \(X=C/H'\), and let
\(P_\chi\) be the primitive cyclotomic factor of \(J(C_\chi)\), with
\(F=\mathbb Q(\zeta_m)\).  If its \(F\)-rank \(h\) in \(H^1\) is
positive, then \(h\) is even, each primitive embedding has Hodge numbers
\((h/2,h/2)\), and \(\bigwedge_F^hH^1(P_\chi,\mathbb Q)\) is algebraic.

\begin{proof}
The morphism \(C_\chi\to X\) is a connected cyclic \(m\)-cover.  Its
inertia groups are images of intersections with conjugates of the original
order-two inertia groups.  Thus every nontrivial image is \(-1\in\mu_m\).
For odd \(m\) the cover is unramified.  For even \(m\), every nonzero local
exponent is \(m/2\).  The relation in the abelianized punctured surface
group says that the sum of the exponents is zero modulo \(m\), so their
number \(s\) is even.  They therefore pair with their negatives.

At every primitive embedding, the nonzero eigenangles are all \(1/2\).
Writing \(q=g(X)\), cyclic-cover Riemann--Roch gives
\[
                     h=2q-2+s,\qquad
       h^{1,0}_\chi=h^{0,1}_\chi=q-1+\frac{s}{2}=\frac h2.
\]
Thus the cover is simple in Schoen's exact definition on p.~11 of
\cite{G34:Schoen}.  His Theorem 2.0 constructs the determinant classes on
\(C_\chi^h\).  D-G34-04 supplies its exceptional-cohomology step for the
present exponents.  The Jacobian correspondence and the cyclotomic idempotent
then give the determinant subspace on \(P_\chi\), as in his Corollary 3.1,
pp.~24--25.  A zero-dimensional factor is omitted.  All primitive embeddings
have been covered, not only the distinguished character \(\chi\).
\end{proof}

\subsection*{D-G34-06. Brauer induction and the algebraic contraction}

\textbf{Statement.} Both Weil lines of the \(G_{34}\) eightfold are algebraic.
More generally the transfer below applies to a non-self-dual multiplicity
determinant whenever the required factor is realized by a rational
group-algebra idempotent as in D-G34-01.  We make no claim here that every
abstract irreducible character, regardless of Schur index, has such a
one-copy rational factor.

\begin{proof}
Let \(J=J(C)\), \(H=H^1(J,\mathbb Q)\), and
\(U=\operatorname{Cent}_{\operatorname{Sp}(H)}(G)\), the full algebraic
centralizer.  Decompose \(H_{\mathbb C}\) into irreducible \(G\)-types.
A non-self-dual pair \(\tau,\overline\tau\) has multiplicity spaces
\(M_\tau,M_\tau^*\), giving a general linear factor of \(U_{\mathbb C}\).
A self-dual orthogonal \(G\)-type gives a symplectic multiplicity factor;
a self-dual quaternionic type \(\eta\) gives an orthogonal multiplicity
factor.  Only types occurring in \(H\) need be retained.  Consequently
\[
 X^*(U_{\mathbb C})=
       \mathbb Z^{\{\text{non-self-dual pairs}\}}
       \oplus(\mathbb Z/2)^{\{\text{quaternionic types}\}}.
\]
The disconnected orthogonal factors must be included.

For a complex \(G\)-representation \(V\), define \(a(V)\) by the free
coordinates
\(\operatorname{mult}_\tau(V)-\operatorname{mult}_{\overline\tau}(V)\)
and torsion coordinates \(\operatorname{mult}_\eta(V)\bmod2\).
This is additive on the representation ring.  Frobenius reciprocity identifies
\(\operatorname{Hom}_{H'}(\chi,H_{\mathbb C})\) with the
\(\chi\)-eigenspace of \(H^1(C_\chi,\mathbb C)\).  The character of its
determinant under \(U\) is precisely \(a(\operatorname{Ind}_{H'}^G\chi)\):
the general linear factors contribute the indicated determinant powers,
the symplectic factors have no character, and the orthogonal factors
contribute the indicated parities.

Brauer's integral induction theorem
\cite[\S10.5, Theorems 19--20]{G34:Serre} expresses the natural irreducible
character as
\[
                          \sigma=\sum_i n_i
                                  \operatorname{Ind}_{H_i}^G\chi_i,
                             \qquad n_i\in\mathbb Z.
\]
Applying \(a\) gives
\[
                           e_\sigma=\sum_i n_i
                             a(\operatorname{Ind}_{H_i}^G\chi_i),
\]
where \(e_\sigma\) is the unit determinant character for \(M_\sigma\).
Characters of orders one and two contribute zero.  Their free coordinates
are zero because they are real.  On a quaternionic irreducible, their
eigenspaces inherit a nondegenerate alternating form: distinct eigencharacters
are orthogonal unless their product is one, and a real eigencharacter is
its own inverse.  Their dimensions are therefore even, giving zero parity
coordinates.  A term whose eigenspace in \(H_{\mathbb C}\) is zero also
contributes zero.  Every remaining term is supplied by D-G34-05.
Replacing \(\chi_i\) by \(\chi_i^{-1}\) negates its full character:
the free coordinates change sign, and the parity coordinates are their own
negatives.  This deals with every negative coefficient.

We spell out the contraction.  For each retained supplier choose a nonzero
rational algebraic class \(y_i\) in its cyclotomic determinant space.
Its components at all primitive embeddings are nonzero.  Indeed these
components are the embeddings of a rank-one \(F_i\)-space, and a nonzero
rational element has no zero embedding.  Choose the embedding determinant
line \(L_i\subset H^{h_i}(P_i,\mathbb C)\) with character
\(\operatorname{sign}(n_i)a(\operatorname{Ind}\chi_i)\).
The complementary top exterior line
\[
                     L_i^\vee\subset
                       H^{2\dim P_i-h_i}(P_i,\mathbb C)
\]
has the inverse character and pairs perfectly with \(L_i\).  This follows
also by taking the product of all other embedding determinant lines: the
full determinant of \(H^1(P_i)\) is trivial under \(U\), because \(U\)
preserves the polarization on the factor.

On \(B=\prod_iP_i^{|n_i|}\), take a pure tensor \(z\) in these complementary
lines, and let \(x\ne0\) lie in \(\det M_\sigma\subset H^8(A,\mathbb C)\).
The character identity makes \(\Gamma=z\otimes x\) on \(B\times A\)
\(U\)-invariant.  Each supplier is an algebraic isogeny factor of \(J(C)\):
use quotient-curve pullback and its rational cyclotomic projector.
The factor \(A\) is similarly given by a rational group-algebra idempotent.
These correspondences commute with \(U\).  A rational \(U\)-equivariant
splitting on \(H^1\), followed by exterior powers, lifts \(\Gamma\) to a
\(U\)-invariant class on a power of \(J\).

The Lefschetz group \(S(J)\) is contained in \(U\), since
\(\operatorname{End}^0(J)\) contains \(\mathbb Q[G]\).
Milne's theorem \cite[Theorem 3.2; Lemma 3.1]{G34:Milne} says that the
\(S(J)\)-invariants on all powers of \(J\) are generated by divisor
classes, also after extension to complex coefficients.  The lifted class,
and hence \(\Gamma\), is therefore algebraic with complex coefficients.
Contract it with the rational algebraic class
\(\boxtimes_i y_i^{\boxtimes|n_i|}\) and push forward along \(B\).
Only the chosen primitive component pairs with each complementary line.
The result is a nonzero scalar multiple of \(x\), so \(x\) is algebraic
with complex coefficients.  Repeat for its conjugate line.  Both the rational
Weil space and the rational cycle-class image are defined over \(\mathbb Q\);
inclusion after extension to \(\mathbb C\) implies inclusion over
\(\mathbb Q\).  This proves algebraicity of the whole Weil plane.

For the quantitative construction there is a useful middle-degree version.
On the curve power \(C_i^{h_i}\), the determinant line is contained in
\(H^1(C_i)^{\otimes h_i}\).  The determinant line of the conjugate
eigenspace has inverse character and is its dual under integration.  Both
have degree \(h_i\).  Thus, with
\(\ell=\sum_i|n_i|h_i\), the contracting tensor on
\((\prod_i C_i^{h_i|n_i|})\times A\) has degree \(\ell+8\).
It lifts to a \(U\)-invariant \(H^1\)-tensor on a power of \(J\) by the
same correspondences.  This is the precise interface to the supplied degree
estimate; no inverse Lefschetz operation with uncontrolled denominators
is needed.
\end{proof}

\subsection*{D-G34-07. An explicit effective bound and the order correction}

The full proofs of \cite[SQ1--SQ7]{G34:Hecke} have been read and checked in
this run by two readers.  Their geometric premises are now supplied by
D-G34-01--06, with no Severn theorem premise: a split rank-eight factor,
rank-sixteen dominance, algebraic cyclic determinant suppliers with every
primitive component nonzero, and the integral relation in the full
centralizer character group.  The arithmetic and intersection estimates of
SQ1--SQ7 are used as supplied results, with the following explicit definitions.

Set
\[
\begin{gathered}
 G=39\,191\,040,\qquad g=4G+1,\qquad
 R=(2G+1)G!G^G,\qquad L=2gR,\\
 S_0=G^{2g}(g!)^2,\qquad S_B=(L/2)!S_0^R,\qquad k=L/2+4,\\
 B_0=(40G)^k\binom{k}{4}S_B\,8!G^8.
\end{gathered}
\]
If the target lattice is \(\mathcal O\)-stable and its polarization has
elementary divisors \(d_1,\ldots,d_8\), write
\(\Delta=\prod_i d_i\) and define
\[
             N=4G^8\Delta,\qquad
             \boxed{D_8(\Delta)=N^8B_0.}
\]
These are positive integers.  The precise estimates used are: coefficient
mass at most \(R\) in the integral character relation (SQ3); supplier ranks
\(h_i\le g\) and integral orbit-sum/Ramanujan-projector mass at most
\(S_0\) (SQ4); exterior-product mass at most \(S_B\) (SQ5); globally
generated divisor representatives bounded numerically by \(40G\) times the
product polarization and the corresponding mixed-degree estimate \(B_0\)
(SQ6); and a split-lattice comparison cleared by \(4G^8\Delta\), using
\(\Delta_A\mid G^8\) for the source (SQ2, SQ7).

Consequently a dominant analytic open subset of the fixed-data Weil family
has an effective integral four-cycle of polarization degree at most
\(D_8(\Delta)\) with nonzero Weil projection.  The accepted bounded-cycle
closedness and transport results of the project extend the same bound over
the connected fixed-data family.  The accepted Theorem W then gives, in
its normalization of denominator \(\delta\) and norm bound \(T\),
\[
       \delta_W\le D_8(\Delta),\qquad
       c_t\le \left\lceil C_{4,3+\omega}(\delta,T)\right\rceil
                       D_8(\Delta)^4.
\]
If the chosen definition of \(C\) is already integral, the ceiling is
redundant.  No bound here is obtained merely from algebraicity point by point;
it is the uniform supplier/intersection construction and the proved
rank-sixteen source that give the degree bound.

\paragraph{D-G34-07b: a nonmaximal endomorphism order.}
Suppose the target lattice \(\Lambda\) is stable under
\(\mathcal O_f=\mathbb Z+f\mathcal O\).  Define
\(\Lambda'=f\mathcal O\Lambda\).  Then
\[
                  f\Lambda\subset\Lambda'\subset\Lambda,
      \qquad I=[\Lambda:\Lambda']\le f^{16},
\]
and \(\Lambda'\) is \(\mathcal O\)-stable.  The identity on the rational
vector space induces an isogeny
\(\phi:A'=V_{\mathbb R}/\Lambda'\to A\).
The pulled-back polarization \(\theta'=\phi^*\theta\) is integral,
positive, and compatible with \(\mathcal O\).  Its elementary-divisor
product is
\(\Delta'=I\Delta\), because its polarization matrix is the restriction
of the original alternating form to an index-\(I\) sublattice.  Thus
\(\Delta'\le f^{16}\Delta\).
Apply the maximal-order bound to \(A'\) and push its effective integral
four-cycle forward by \(\phi\).  The projection formula gives
\[
                 \deg_\theta(\phi_*Z)=\deg_{\theta'}Z.
\]
An isogeny induces an isomorphism on rational cohomology and carries the
Weil subspace to the Weil subspace, so the nonzero Weil projection remains
nonzero.  Since \(D_8\) is increasing in its positive integer argument,
\[
 \boxed{\delta_W(A,\theta)\le D_8(f^{16}\Delta),\qquad
 c_t\le\left\lceil C_{4,3+\omega}(\delta,T)\right\rceil
                         D_8(f^{16}\Delta)^4.}
\]
This proves the fixed-order extension without suppressing a conductor factor.
A rational \(K\)-action on a fixed integral local system has some such fixed
order; the bound is uniform with that order and polarization data fixed.
